\chapter{Firmware}\label{ch-firmware}

Firmware is a program that runs on a board with no operating system
under it: it starts when the power comes on and runs until it goes
off. It is held to rules a desktop program is not, and often to a
written coding standard such as MISRA C. \verb|zeno build --firmware|
writes C for that world. This chapter is about what it changes, how to
put a real board under the program, and what it does not yet claim.

\section{Errors as values}

An ordinary build compiles \verb|try| with C's \verb|setjmp| and
\verb|longjmp|: a \verb|raise| jumps straight back to the \verb|try|,
out of however many functions are in between. That is cheap, and it is
exactly what standards for safety-critical firmware forbid. MISRA C
bans \verb|<setjmp.h>| outright, because the jump skips what the
functions in between still meant to do, can leave local variables with
no defined value, and hides paths from the tools that measure stack
depth and test coverage.

A firmware build carries errors as ordinary return values instead. The
Zeno does not change:

\begin{zeno}[file=speed.zeno, cmd=zeno build --firmware {file} ;; ./speed ;; grep setjmp speed.c | wc -l]
function parse_speed(text: str) -> int raises ValueError {
    let n = int(text)
    if n < 0 or n > 100 {
        raise ValueError(f"a speed is 0 to 100, and {n} is not")
    }
    return n
}

for text in ["40", "fast", "250"] {
    try {
        print("speed", parse_speed(text))
    } except ValueError as e {
        print("refused:", e.message)
    }
}
\end{zeno}

It prints what \verb|zeno run| prints, and there is no \verb|setjmp|
in the C (\verb|grep| found 0 lines). Here is what \verb|parse_speed|
became:

\begin{zeno}[file=speed.zeno, show=no, cmd=zeno build --firmware {file} > /dev/null ;; grep -A11 "^static bool z_parse_speed.*{$" speed.c]
function parse_speed(text: str) -> int raises ValueError {
    let n = int(text)
    if n < 0 or n > 100 {
        raise ValueError(f"a speed is 0 to 100, and {n} is not")
    }
    return n
}

for text in ["40", "fast", "250"] {
    try {
        print("speed", parse_speed(text))
    } except ValueError as e {
        print("refused:", e.message)
    }
}
\end{zeno}

Three things changed. The function returns \verb|true| when it worked
and \verb|false| when it raised, and its answer goes out through one
more parameter. A \verb|raise| notes the error, with its line and
column, and returns \verb|false|. And a call to something that can
raise, \verb|int(text)| here, is followed at once by a test. Only the
functions that say \verb|raises| are compiled this way; the rest are
plain C, with no tests after calls to them. That is why \verb|raises|
must be written in a build (\hyperref[ch-errors]{When Things Go Wrong}).

\section{Faults}

Zeno has two sorts of error that \verb|zeno run| can catch:

\begin{itemize}
\item errors \textbf{the program raises itself}: \verb|raise|, and
  \verb|int| or \verb|float| of text;
\item errors \textbf{the runtime finds}: a division by zero, a position
  past the end, a key that is not there, an overflow, \verb|pop| on an
  empty list.
\end{itemize}

A firmware build carries the first sort as values. The second sort is
a \emph{fault}: the program stops there, with the message \verb|zeno run|
gives for an uncaught error, and no \verb|except| sees it. To carry a
division by zero as a value, every function would have to be one that
can fail, and nearly every line would need a test after it. The
firmware answer is to test first:

\begin{zenofrag}
if divisor != 0 {
    rate = total // divisor
}
\end{zenofrag}

The compiler does not leave this to be found when the board resets. A
\verb|try| that hopes to catch a fault is refused before anything is
built:

\begin{zeno}[file=divide.zeno, cmd=zeno build --firmware {file}, fails=yes]
let readings = [4, 0]
for n in readings {
    try {
        print(100 // n)
    } except ZeroDivisionError {
        print("no reading")
    }
}
\end{zeno}

A \verb|FileError| is the exception to the exception: no check the
program could make rules out a missing file, so it is always a value,
and can always be caught.

\begin{why}
The line between the two sorts is the line between what the program
can check beforehand and what it cannot. A bad index can be ruled out
by testing the length; a missing file cannot. Firmware held to a
standard is expected to test, and the build makes the cases it did not
test visible. (Decisions 4.1 and 5.5l.)
\end{why}

\section{A real board underneath}

The board module meets the hardware in four C functions
(\hyperref[ch-board]{The Board}). For a real board, write them in a
header over the vendor's library, and tell the C compiler where it is
with \verb|-DZENO_BOARD|:

\begin{zenofrag}[lang=shell]
$ zeno build --firmware --no-collector blink.zeno
$ cc -O2 -DZENO_BOARD='"my_board.h"' -o blink blink.c -lm
\end{zenofrag}

With ST's HAL for an STM32, where the LED of a Nucleo board is on port
A pin 5, the header would look like this. It is a sketch, and has not
been run on hardware:

\begin{zenofrag}[lang=shell]
#define ZENO_PINS 16

static void zeno_board_write(int pin, bool level) {
    if (pin == ZENO_LED) pin = 5;
    HAL_GPIO_WritePin(GPIOA, (uint16_t)(1u << pin),
                      level ? GPIO_PIN_SET : GPIO_PIN_RESET);
}
static bool zeno_board_read(int pin) {
    if (pin == ZENO_LED) pin = 5;
    return HAL_GPIO_ReadPin(GPIOA, (uint16_t)(1u << pin)) == GPIO_PIN_SET;
}
static void zeno_board_sleep(double seconds) {
    HAL_Delay((uint32_t)(seconds * 1000.0));
}
static double zeno_board_now(void) { return (double)HAL_GetTick(); }
\end{zenofrag}

To the Zeno program the LED is always pin 13, and the board is free to
wire that number to wherever its LED really is. The same flag works for
an LLVM build, on the \verb|clang| line.

A board whose clock moves only when the program sleeps makes a program
that waits run in no time. That is how firmware is tested on a desk,
and Zeno's own tests do exactly that.

The same header can take the place of the C library's streams. With
\verb|#define ZENO_BOARD_PRINTS| it brings four more functions, and
\verb|print|, \verb|input|, a warning, the message of an error and the
end of the program all go through them. For a board with one serial
port and nothing to read from:

\begin{zenofrag}[lang=shell]
#define ZENO_BOARD_PRINTS

static void zeno_board_print(const char *bytes, size_t count) {
    HAL_UART_Transmit(&huart2, (uint8_t *)bytes, (uint16_t)count, 100);
}
static void zeno_board_warn(const char *bytes, size_t count) {
    zeno_board_print(bytes, count);
}
static int zeno_board_typed(void) { return -1; }
static void zeno_board_stop(int status) {
    (void)status;
    NVIC_SystemReset();
}
\end{zenofrag}

This too is a sketch for an STM32, and has not been run on hardware.
\verb|boards/lm3s6965evb| in Zeno's repository is a whole header of
this kind, on the registers of the chip QEMU emulates, and Zeno's
tests run it there.

\section{What the C of blink looks like}

All of \verb|blink|'s program, in a firmware build with no collector,
is three lines of C:

\begin{zenofrag}[lang=shell]
while (true) {
    zeno_pin_toggle(ZENO_LED);
    zeno_sleep((500.0 / 1000.0));
}
\end{zenofrag}

A pin is a number, \verb|500 ms| has become seconds, and nothing in it
takes memory or can raise.

\section{What this does not claim}

Removing \verb|setjmp| removes one obstacle to MISRA, not all of them.
The C that a firmware build writes still uses variable
argument lists in the runtime, calls \verb|snprintf| for a number with
a fraction (and \verb|printf| and \verb|exit| on a board that does not
print by itself), and
uses GCC and Clang built-ins for checked arithmetic. Taking memory from
the heap, which MISRA also bans, is answered by \verb|--no-collector|,
in the next chapter. None of these is hidden, and each is a later piece
of work. A function that can raise cannot yet be passed around as a
value in a firmware build, because a function type does not say whether
it raises.

\begin{summary}
\verb|zeno build --firmware| carries errors the program raises as
return values, tested after each call, with no \verb|setjmp|. Errors
the runtime finds are faults that stop the program, and a \verb|try|
that hopes to catch one is refused. A real board goes under the program
as a header of four functions, named with \verb|-DZENO_BOARD|, or by
the \verb|header| line of the board's folder when
\verb|zeno build --board| does the whole build. Four more, behind
\verb|ZENO_BOARD_PRINTS|, let the board print, read and stop by
itself, with no C library stream under the program.
\end{summary}

\begin{tryit}
\begin{enumerate}
\item Build \verb|speed.zeno| with and without \verb|--firmware|, and
  compare the C of \verb|parse_speed| in the two.
\item Rewrite the \verb|divide.zeno| example so that a firmware build
  accepts it.
\item Write a board header of your own for the simulated board that
  prints \verb|HIGH| and \verb|LOW|, and build \verb|blink| with it.
\end{enumerate}
\end{tryit}
